% GENERATED FILE. Edit canonical lessons, then run npm run generate:books.
% canonical-source-hash: fnv1a64:0acf5a13b63f86c5
% canonical-lessons: MW-C241-chot, MW-C241-thokar, MW-C241-ghav, MW-C241-tutno, MW-C241-khisakno

\chapter{A Fall in the Lane}
\label{ch:mw-a-fall-in-the-lane}
\hlchaptermodality{241}

\begin{chapteropening}
\textbf{By the end of this chapter:} \emph{I can say that I fell and got hurt.}

Complete the last lesson of chapter 241: I can say that I fell and got hurt.
\end{chapteropening}

\section[\texorpdfstring{\mw{चोट} (choṭ)}{चोट (choṭ)}]{\mw{चोट} (choṭ) — an injury; a hurt}

\label{lesson:MW-C241-chot}



\begin{quote}
Before the new one: say the Marwadi for brave, then the Marwadi for crazy.
\end{quote}

\subsection*{You'll want to know: \mw{चोट}}
\textbf{\mw{चोट}} — \emph{choṭ} — \textquotedblleft{}an injury; a hurt\textquotedblright{}.

\begin{quote}
\mw{म्हारी} \mw{चोट} \mw{घणी} \mw{माड़ी} \mw{कोनी।}
\end{quote}

\begin{quote}
\emph{mhārī choṭ ghaṇī māṛī konī.}
\end{quote}

\begin{quote}
My injury is not very bad.
\end{quote}

A feminine noun.

\begin{cousinweb}
\textbf{\mw{चोट}} (\emph{choṭ}) is inherited from Sanskrit \emph{cuṭati}, \textquotedblleft{}cuts off, strikes\textquotedblright{}. CDIAL gives no Marwari form in this entry; it lists Hindi \emph{coṭ}, \textquotedblleft{}wound\textquotedblright{} (CDIAL 4857).
\end{cousinweb}

\begin{grammarlens}[title={What you've built}]
One more word for saying what happened.
\end{grammarlens}

\subsection*{Your turn}
\begin{itemize}
\raggedright
  \item \emph{Say it:} \emph{choṭ}
  \item \emph{Say it:} \emph{choṭ}, once more
  \item \emph{Say it:} \emph{mhārī choṭ ghaṇī māṛī konī.} — My injury is not very bad.
  \item \emph{Recall:} say the Marwadi for brave, then the Marwadi for crazy, then say \emph{choṭ} again
\end{itemize}

\begin{tcolorbox}[breakable,colback=teal!4,colframe=teal!35!black,title={Before you move on}]
What does \mw{चोट} mean? (\textquotedblleft{}An injury; a hurt\textquotedblright{}.) Say it once more.
\end{tcolorbox}
\section[\texorpdfstring{\mw{ठोकर} (ṭhokar)}{ठोकर (ṭhokar)}]{\mw{ठोकर} (ṭhokar) — a stumble; a knock against the foot}

\label{lesson:MW-C241-thokar}



\begin{quote}
Before the new one: say the Marwadi for crazy, then the Marwadi for an injury.
\end{quote}

\subsection*{You'll want to know: \mw{ठोकर}}
\textbf{\mw{ठोकर}} — \emph{ṭhokar} — \textquotedblleft{}a stumble; a knock against the foot\textquotedblright{}.

\begin{quote}
\mw{काल} \mw{सड़क} \mw{माथै} \mw{ठोकर} \mw{लागी।}
\end{quote}

\begin{quote}
\emph{kāl saṛak māthai ṭhokar lāgī.}
\end{quote}

\begin{quote}
I stumbled on the road yesterday.
\end{quote}

A feminine noun.

\begin{cousinweb}
\textbf{\mw{ठोकर}} (\emph{ṭhokar}) is inherited from a reconstructed \emph{*ṭhōkk}, \textquotedblleft{}knock\textquotedblright{}. CDIAL gives no Marwari form in this entry; it lists Hindi \emph{ṭhokar}, \textquotedblleft{}blow\textquotedblright{} (CDIAL 5513).
\end{cousinweb}

\begin{grammarlens}[title={What you've built}]
One more word for saying what happened.
\end{grammarlens}

\subsection*{Your turn}
\begin{itemize}
\raggedright
  \item \emph{Say it:} \emph{ṭhokar}
  \item \emph{Say it:} \emph{ṭhokar}, once more
  \item \emph{Say it:} \emph{kāl saṛak māthai ṭhokar lāgī.} — I stumbled on the road yesterday.
  \item \emph{Recall:} say the Marwadi for crazy, then the Marwadi for an injury, then say \emph{ṭhokar} again
\end{itemize}

\begin{tcolorbox}[breakable,colback=teal!4,colframe=teal!35!black,title={Before you move on}]
What does \mw{ठोकर} mean? (\textquotedblleft{}A stumble; a knock against the foot\textquotedblright{}.) Say it once more.
\end{tcolorbox}
\section[\texorpdfstring{\mw{घाव} (ghāv)}{घाव (ghāv)}]{\mw{घाव} (ghāv) — a wound}

\label{lesson:MW-C241-ghav}



\begin{quote}
Before the new one: say the Marwadi for an injury, then the Marwadi for a stumble.
\end{quote}

\subsection*{You'll want to know: \mw{घाव}}
\textbf{\mw{घाव}} — \emph{ghāv} — \textquotedblleft{}a wound\textquotedblright{}.

\begin{quote}
\mw{घाव} \mw{छोटो} \mw{है}, \mw{पण} \mw{दरद} \mw{घणो} \mw{है।}
\end{quote}

\begin{quote}
\emph{ghāv chhoṭo hai, paṇ darad ghaṇo hai.}
\end{quote}

\begin{quote}
The wound is small, but it hurts a lot.
\end{quote}

A masculine noun.

\begin{cousinweb}
\textbf{\mw{घाव}} (\emph{ghāv}) is inherited from Sanskrit \emph{ghāta}, \textquotedblleft{}a blow, a bruise\textquotedblright{}. CDIAL gives no Marwari form in this entry; it lists Hindi \emph{ghāu}, \textquotedblleft{}wound\textquotedblright{} (CDIAL 4460).
\end{cousinweb}

\begin{grammarlens}[title={What you've built}]
One more word for saying what happened.
\end{grammarlens}

\subsection*{Your turn}
\begin{itemize}
\raggedright
  \item \emph{Say it:} \emph{ghāv}
  \item \emph{Say it:} \emph{ghāv}, once more
  \item \emph{Say it:} \emph{ghāv chhoṭo hai, paṇ darad ghaṇo hai.} — The wound is small, but it hurts a lot.
  \item \emph{Recall:} say the Marwadi for an injury, then the Marwadi for a stumble, then say \emph{ghāv} again
\end{itemize}

\begin{tcolorbox}[breakable,colback=teal!4,colframe=teal!35!black,title={Before you move on}]
What does \mw{घाव} mean? (\textquotedblleft{}A wound\textquotedblright{}.) Say it once more.
\end{tcolorbox}
\section[\texorpdfstring{\mw{टूटणो} (ṭūṭṇo)}{टूटणो (ṭūṭṇo)}]{\mw{टूटणो} (ṭūṭṇo) — to get broken; to break by itself}

\label{lesson:MW-C241-tutno}



\begin{quote}
Before the new one: say the Marwadi for a stumble, then the Marwadi for a wound.
\end{quote}

\subsection*{You'll want to know: \mw{टूटणो}}
\textbf{\mw{टूटणो}} — \emph{ṭūṭṇo} — \textquotedblleft{}to get broken; to break by itself\textquotedblright{}.

\begin{quote}
\mw{म्हारी} \mw{चूड़ी} \mw{टूटी।}
\end{quote}

\begin{quote}
\emph{mhārī chūṛī ṭūṭī.}
\end{quote}

\begin{quote}
My bangle broke.
\end{quote}

Past (it …): \mw{वो} \mw{प्यालो} \mw{टूट्यो}, \emph{vo pyālo ṭūṭyo}.

\begin{cousinweb}
\textbf{\mw{टूटणो}} (\emph{ṭūṭṇo}) is inherited from Sanskrit \emph{truṭyati}, \textquotedblleft{}is broken\textquotedblright{}; CDIAL lists Old Marwari \emph{tūṭaï}, \textquotedblleft{}is broken\textquotedblright{}, under it (CDIAL 6065).
\end{cousinweb}

\begin{grammarlens}[title={What you've built}]
One more word for saying what happened.
\end{grammarlens}

\subsection*{Your turn}
\begin{itemize}
\raggedright
  \item \emph{Say it:} \emph{ṭūṭṇo}
  \item \emph{Say it:} \emph{ṭūṭṇo}, once more
  \item \emph{Say it:} \emph{mhārī chūṛī ṭūṭī.} — My bangle broke.
  \item \emph{Recall:} say the Marwadi for a stumble, then the Marwadi for a wound, then say \emph{ṭūṭṇo} again
\end{itemize}

\begin{tcolorbox}[breakable,colback=teal!4,colframe=teal!35!black,title={Before you move on}]
What does \mw{टूटणो} mean? (\textquotedblleft{}To get broken; to break by itself\textquotedblright{}.) Say it once more.
\end{tcolorbox}
\section[\texorpdfstring{\mw{खिसकणो} (khisakṇo)}{खिसकणो (khisakṇo)}]{\mw{खिसकणो} (khisakṇo) — to slip; to slide out of place}

\label{lesson:MW-C241-khisakno}



\begin{quote}
Before the new one: say the Marwadi for a wound, then the Marwadi for to get broken.
\end{quote}

\subsection*{You'll want to know: \mw{खिसकणो}}
\textbf{\mw{खिसकणो}} — \emph{khisakṇo} — \textquotedblleft{}to slip; to slide out of place\textquotedblright{}.

\begin{quote}
\mw{सड़क} \mw{माथै} \mw{म्हारो} \mw{पग} \mw{खिसक्यो।}
\end{quote}

\begin{quote}
\emph{saṛak māthai mhāro pag khisakyo.}
\end{quote}

\begin{quote}
My foot slipped on the road.
\end{quote}

Past (I …): \mw{हूं} \mw{खिसक्यो}, \emph{hū̃ khisakyo}.

\begin{cousinweb}
\textbf{\mw{खिसकणो}} (\emph{khisakṇo}) is inherited from a reconstructed \emph{*khis}, \textquotedblleft{}slip\textquotedblright{}. CDIAL gives no Marwari form in this entry; it lists Hindi \emph{khisaknā}, \textquotedblleft{}to slip away\textquotedblright{} (CDIAL 3888).
\end{cousinweb}

\begin{grammarlens}[title={What you've built}]
One more word for saying what happened.
\end{grammarlens}

\subsection*{Your turn}
\begin{itemize}
\raggedright
  \item \emph{Say it:} \emph{khisakṇo}
  \item \emph{Say it:} \emph{khisakṇo}, once more
  \item \emph{Say it:} \emph{saṛak māthai mhāro pag khisakyo.} — My foot slipped on the road.
  \item \emph{Recall:} say the Marwadi for a wound, then the Marwadi for to get broken, then say \emph{khisakṇo} again
\end{itemize}

\begin{tcolorbox}[breakable,colback=teal!4,colframe=teal!35!black,title={Before you move on}]
What does \mw{खिसकणो} mean? (\textquotedblleft{}To slip; to slide out of place\textquotedblright{}.) Say it once more.
\end{tcolorbox}
